\chapter*{Abstract}

Today's computer component market is developing very dynamically, which leads to an increasing number of their types, models or configurations.
Assembling a personal computer might seem like a complicated task, considering hardware compatibility and differences in the individual components construction.
Learning to assemble a computer on real hardware also carries the risk of damaging expensive components, especially when handling delicate parts such as the processor.
Traditional learning methods, such as books and videos, do not provide the necessary knowledge and practical experience.
This leads to a growing demand for an interactive learning environment that allows users to learn how to build a computer without bearing financial risks.
An application using virtual reality technology, introducing users to the topic of computer assembly, may be the answer to these problems.

The objective of the diploma project "Virtual assembly and visualization of the operation of a personal computer" was to design and implement a virtual trainer operating in the virtual reality caves located at the Immersive 3D Visualization Lab at the Gdańsk University of Technology.
As a result of the work, a computer building simulator was created, offering its users the ability to virtually assemble multiple hardware configurations, taking into account component compatibility, and guiding them through the entire process in a realistic way.
The application provides a safe environment for learning how to build a computer, without bearing the risk of damaging expensive components, and the ability to freely repeat completed steps reinforces acquired knowledge.
The use of virtual reality technology allowed for the accurate detection and mapping of the user's hand and head positions, providing a high level of immersion in the created environment.
The trainer was implemented with the necessary mechanics for full control over computer components, allowing them to be freely moved and connected.
To provide user support, application interface elements were created that guide users through each stage of building a virtual computer in an accessible manner.
The simulator was created using the Unreal Engine 5 game engine, which allows for direct interaction with motion tracking systems and multi-screen projection synchronization.

The simulator was thoroughly tested in virtual reality caves, and a research was conducted to assess its intuitiveness and educational value.
User tests demonstrated high levels of satisfaction, confirming the trainer's effectiveness as an educational tool.
Additionally, comparative research was conducted between the simulator and a commercially available solution, which showed similar or higher levels of user satisfaction with the virtual reality application.
Participants positively assessed the trainer's intuitiveness and demonstrated greater confidence in assembling the computer themselves.

The diploma project was realized by a team of four students.
Each team member was assigned the following tasks during the implementation of the trainer and the diploma thesis completion:
\begin{itemize}
    \item \textbf{Jakub Jędrzejczyk} - was responsible for team management and implementing the mechanics for controlling components and assembling the computer, as well as creating a legible user interface. He authored chapter \ref{chap:Testing}, and sections \ref{sec:GameEngineSelection}, \ref{chap:GUI}, \ref{chap:ObjectHolding}, \ref{chap:Dtrack}, \ref{chap:ConnectingComponents}, \ref{chap:PCValidation}, \ref{sec:FuturePlans}.
    \item \textbf{Jakub Kinder} - was responsible for designing and implementing a build manager, which tracked the progress of the computer assembly and validated the built configuration. He authored chapters \ref{chap:Introduction} and \ref{chap:FieldIntroduction}, and sections \ref{chap:buildmanager}. 
    \item \textbf{Hubert Leszczyński} - was responsible for a system project, designing the main scene of the application, and creating the 3D graphics. He also implemented the mechanic of a dynamic computer component creation. He authored chapter \ref{chap:project} and sections \ref{chap:models}, \ref{chap:spawners}, \ref{chap:objectConfig} and \ref{chap:sceneDesign}.
    \item \textbf{Mateusz Merda} - was responsible for developing objects with advanced physics and collision configurations. He authored chapter 3 and sections \ref{subsec:twohands}, \ref{sec:cableconnector}, \ref{sec:results}, \ref{sec:conclusions}.
\end{itemize}

In the future, the project can be expanded to include precise mechanisms of visualizing the work of a personal computer, as well as a larger number of available hardware components and configurations.
\newline Keywords: Virtual reality, Application development, Visualization, Computer assembly
\newline Field of Science and Technology (OECD classification): 1.2 Computer and information sciences